\section{Полнота. Замкнутые классы. Полные системы функций и базисы. Теорема Поста.}

\subsection*{Функциональная полнота и замкнутость}

\paragraph{Определения}
\begin{itemize}
    \item \textbf{Замыкание} $[F]$ системы функций $F$ — это множество всех функций, которые можно реализовать формулами в базисе $F$.
    \item \textbf{Замкнутый класс} — это множество функций $K$, которое совпадает со своим замыканием ($[K] = K$). То есть суперпозиция функций из этого класса не выводит за его пределы.
    \item \textbf{Полная система} — это система функций $F$, замыкание которой совпадает со всем множеством булевых функций ($[F] = P_n$).
    \item \textbf{Базис} — это полная система, которая перестает быть полной при удалении любого её элемента.
\end{itemize}

\subsection*{Пять предполных классов Поста}

Новиков выделяет пять классов функций, обладающих свойством замкнутости:
\begin{enumerate}
    \item $T_0$ — \textbf{Класс функций, сохраняющих нуль:} $f(0, \dots, 0) = 0$.
    \item $T_1$ — \textbf{Класс функций, сохраняющих единицу:} $f(1, \dots, 1) = 1$.
    \item $T^*$ — \textbf{Класс самодвойственных функций:} $f = f^*$.
    \item $T_{\le}$ — \textbf{Класс монотонных функций:} $\alpha \preceq \beta \implies f(\alpha) \le f(\beta)$.
    \item $T_L$ — \textbf{Класс линейных функций:} функции, представимые полиномом Жегалкина степени $\le 1$.
\end{enumerate}

\subsection*{Теорема Поста}

\paragraph{Формулировка}
Система булевых функций $F$ является \textbf{полной} тогда и только тогда, когда она не содержится целиком ни в одном из пяти замкнутых классов $T_0, T_1, T^*, T_{\le}, T_L$.

\paragraph{Доказательство (схема)}
\begin{enumerate}
    \item \textbf{Необходимость:} Если система $F$ целиком лежит в одном из классов (например, все функции сохраняют 0), то и любая их суперпозиция будет сохранять 0. Тогда мы не сможем получить функцию $\neg x$, так как $\neg 0 = 1$.
    \item \textbf{Достаточность:} Если функции системы $F$ «разрушают» свойства всех 5 классов, то из них можно построить базис $\{\neg, \land\}$:
    \begin{itemize}
        \item С помощью функций $\notin T_0$ и $\notin T_1$ строятся \textbf{константы} $0$ и $1$.
        \item С помощью функции $\notin T_{\le}$ (немонотонной) и констант строится \textbf{отрицание}.
        \item С помощью функции $\notin T_L$ (нелинейной), отрицания и констант строится \textbf{конъюнкция}.
    \end{itemize}
\end{enumerate}

\subsection*{Примеры полных систем}

\begin{enumerate}
    \item $\{\neg, \land, \lor\}$ — избыточная полная система.
    \item $\{\neg, \land\}$ — стандартный базис.
    \item $\{\neg, \lor\}$ — стандартный базис.
    \item $\{| \}$ (Штрих Шеффера) — полная система из одной функции.
    \item $\{\downarrow \}$ (Стрелка Пирса) — полная система из одной функции.
    \item $\{x \land y, x \oplus y, 1\}$ — базис Жегалкина.
\end{enumerate}

\paragraph{Принцип двойственности в полноте}
Если система $F$ полна, то система двойственных ей функций $F^*$ также является полной. Это следует из того, что любая функция $h$ может быть выражена через $F$, а значит $h^*$ выразима через $F^*$.

\paragraph{Пример проверки по таблице Поста}
Для системы $\{\neg, \land\}$:
\begin{itemize}
    \item $\neg x$: $\notin T_0$ (т.к. $\neg 0=1$), $\notin T_1$ (т.к. $\neg 1=0$), $\notin T_{\le}$.
    \item $x \land y$: $\in T_0, \in T_1, \in T_{\le}$, но $\notin T^*$ и $\notin T_L$.
    \item \textbf{Итог:} Вместе они перекрывают все 5 классов $\implies$ система полна.
\end{itemize}